% ===== 丘维声《解析几何》讲解 · 统一样式（与第 1–5 批一致） =====
\usepackage[a4paper,top=2.4cm,bottom=2.4cm,left=2.5cm,right=2.5cm,headheight=15pt]{geometry}
\setCJKmainfont{Noto Serif CJK SC}[BoldFont={Noto Serif CJK SC Bold}]
\setCJKsansfont{Noto Sans CJK SC}[BoldFont={Noto Sans CJK SC Bold}]
\setCJKmonofont{Noto Sans Mono CJK SC}
\usepackage{amsmath,amssymb,amsthm,mathtools,bm,mathrsfs}
\usepackage{xcolor}
\usepackage{graphicx}
\usepackage{booktabs,tabularx,array,multirow}
\usepackage{enumitem}
\usepackage{fancyhdr}
\usepackage{titlesec}
\usepackage{caption}
\usepackage{needspace}
\usepackage{tikz}
\usetikzlibrary{calc,intersections,arrows.meta,positioning,decorations.markings,shapes.geometric,through,3d,perspective,patterns}
\usepackage{tikz-3dplot}
\usepackage{pgfplots}\pgfplotsset{compat=1.18}
\usepackage[most]{tcolorbox}
\usepackage[hidelinks,bookmarksnumbered,bookmarksopen]{hyperref}

% ---------- 颜色 ----------
\definecolor{mainblue}{HTML}{1F5FA8}
\definecolor{defblue}{HTML}{2E6DB4}   \definecolor{defbg}{HTML}{EEF3FA}
\definecolor{thmred}{HTML}{B03A2E}    \definecolor{thmbg}{HTML}{FBEFEE}
\definecolor{intor}{HTML}{C8691C}     \definecolor{intbg}{HTML}{FBF4EC}
\definecolor{gdpur}{HTML}{6C4F9E}     \definecolor{gdbg}{HTML}{F4F0F9}
\definecolor{mteal}{HTML}{1E7B6E}     \definecolor{mtbg}{HTML}{EAF5F3}
\definecolor{ckgray}{HTML}{E2E2E2}    \definecolor{ckbg}{HTML}{F7F7F7}
\definecolor{exslate}{HTML}{34495E}   \definecolor{exbg}{HTML}{F4F6F8}
\definecolor{solgreen}{HTML}{2E8B57}
\definecolor{keyred}{HTML}{C0392B}
\definecolor{figblue}{HTML}{2E6DB4}   \definecolor{figred}{HTML}{C0392B}
\definecolor{figgreen}{HTML}{2E8B57}  \definecolor{figor}{HTML}{D9822B}
\definecolor{figpur}{HTML}{7D5BA6}    \definecolor{figgray}{HTML}{7F8C8D}

% ---------- 标题样式 ----------
\ctexset{
  chapter={
    name={}, number={第\arabic{chapter}章},
    format=\centering,
    nameformat=\sffamily\bfseries\Large\color{mainblue},
    titleformat=\sffamily\bfseries\Huge\color{mainblue},
    aftername=\par\vspace{0.8em},
    beforeskip=10pt, afterskip=22pt,
    aftertitle={\par\vspace{6pt}{\color{mainblue}\rule{\textwidth}{0.8pt}}},
  },
  section={
    format=\sffamily\bfseries\Large\color{mainblue}\raggedright,
    name={\S,}, number=\thechapter.\arabic{section},
    aftername=\enspace, beforeskip=18pt, afterskip=10pt,
  },
  subsection={
    format=\sffamily\bfseries\large\color{mainblue!85!black}\raggedright,
    number=\thechapter.\arabic{section}.\arabic{subsection},
    aftername=\enspace, beforeskip=12pt, afterskip=6pt,
  },
}
\setcounter{secnumdepth}{2}
\setcounter{tocdepth}{1}
\renewcommand{\thefigure}{\thechapter-\arabic{figure}}
\captionsetup{labelsep=quad,font=small,labelfont={color=black}}

% 习题详解标题（不编号，进目录与书签）
\newcommand{\xiti}[1]{%
  \vspace{8pt}\phantomsection
  \addcontentsline{toc}{section}{习题 #1 详解}%
  {\noindent\sffamily\bfseries\Large\color{mainblue}习题 #1\ 详解\par}\vspace{8pt}}
\newcommand{\zongjie}[1]{%
  \vspace{8pt}\phantomsection
  \addcontentsline{toc}{section}{#1}%
  {\noindent\sffamily\bfseries\Large\color{mainblue}#1\par}\vspace{8pt}}

% ---------- 页眉页脚 ----------
\pagestyle{fancy}
\fancyhf{}
\fancyhead[L]{\small 高等代数 · 二次型图解}
\fancyhead[R]{\small 第\thechapter 章\quad\leftmark}
\fancyfoot[C]{\small\thepage}
\renewcommand{\headrulewidth}{0.5pt}
\renewcommand{\chaptermark}[1]{\markboth{#1}{}}
\renewcommand{\sectionmark}[1]{}
\fancypagestyle{plain}{\fancyhf{}\fancyfoot[C]{\small\thepage}\renewcommand{\headrulewidth}{0pt}}

% ---------- 彩色框 ----------
\tcbset{
  basebox/.style={enhanced,breakable,arc=2mm,boxrule=0.6pt,
    left=6pt,right=6pt,top=4pt,bottom=4pt,
    fonttitle=\sffamily\bfseries,coltitle=white,
    before skip=8pt,after skip=8pt},
}
\newtcolorbox{defbox}[1]{basebox,colframe=defblue,colback=defbg,colbacktitle=defblue,title={#1}}
\newtcolorbox{thmbox}[1]{basebox,colframe=thmred,colback=thmbg,colbacktitle=thmred,title={#1}}
\newtcolorbox{intbox}[1]{basebox,colframe=intor,colback=intbg,colbacktitle=intor,title={#1}}
\newtcolorbox{gdbox}[1]{basebox,colframe=gdpur,colback=gdbg,colbacktitle=gdpur,title={#1}}
\newtcolorbox{mtbox}[1]{basebox,colframe=mteal,colback=mtbg,colbacktitle=mteal,title={#1}}
\newtcolorbox{ckbox}[1]{basebox,colframe=ckgray,colback=ckbg,colbacktitle=ckgray,coltitle=black,title={#1}}

% 习题题干框：深色题号标签
\newtcolorbox{ex}[1]{enhanced,breakable,arc=2mm,boxrule=0.6pt,
  colframe=exslate,colback=exbg,
  left=6pt,right=6pt,top=6pt,bottom=4pt,
  attach boxed title to top left={xshift=4mm,yshift*=-\tcboxedtitleheight/2},
  boxed title style={colback=exslate,colframe=exslate,sharp corners,boxrule=0pt,left=8pt,right=8pt,top=2pt,bottom=2pt},
  fonttitle=\sffamily\bfseries,coltitle=white,
  title={#1},before skip=14pt,after skip=4pt}

\newcommand{\jie}{\noindent{\color{solgreen}\sffamily\bfseries【解】}}
\newcommand{\zheng}{\noindent{\color{solgreen}\sffamily\bfseries【证】}}
\newcommand{\fenxi}{\noindent{\color{solgreen}\sffamily\bfseries【分析】}}
\newcommand{\gd}{{\color{gdpur}\sffamily\bfseries【高代】}}
\newcommand{\ji}{{\color{intor}\sffamily\bfseries【记】}}
\newcommand{\key}[1]{\leavevmode{\color{keyred}\bfseries #1}}
\renewcommand{\qed}{\hfill$\blacksquare$}

% ---------- 数学记号 ----------
\newcommand{\R}{\mathbb{R}}
\newcommand{\C}{\mathbb{C}}
\newcommand{\pib}{\overline{\pi}_0}
\newcommand{\Sb}{\overline{S}}
\newcommand{\linf}{l_\infty}
\newcommand{\T}{^{\mathsf T}}
\newcommand{\cro}[4]{(#1,#2;#3,#4)}
\newcommand{\vv}[1]{\boldsymbol{#1}}
\newcommand{\pt}[3]{(#1,#2,#3)\T}
\newcommand{\spn}[1]{\langle #1\rangle}
\setlist{nosep,leftmargin=2em}
\setlength{\parindent}{2em}
\linespread{1.25}
